\bookchapter{الفصل 65: المشروع الأول — دفتر حساباتي (نسخة محلية)}

\chapterintro{مرحبًا بك في «مشروعك الجاد» الأول. هذا الدفتر لا يتصل بالشبكة، لكنه مكتمل الأركان: يسجل عملية مالية، ويعرض كشف الحساب، ويجمع الإحصاءات شهريًا، ويتيح الحذف والتعديل. وبإتمامه تكون قد قطعت كامل الدورة «بيانات منطق واجهة». ولا تتسرع في نسخ الكود — فكر بوضوح أولًا ثم انطلق، فهذا مبدؤنا الدائم.}

\booksection{أولًا: التفكير في المتطلبات — اسأل خمسة أسئلة أولًا}

\begin{center}
\begin{tabularx}{\linewidth}{|X|X|}
\hline
\textbf{العنصر} & \textbf{الجواب} \tabularnewline
\hline
Who (لمن يُستخدم) & لنفسك، ويصلح للعائلة والأصدقاء أيضًا \tabularnewline
What (ماذا يفعل) & تسجيل الدخل والمصروف، وعرض كشف الحركة، وجمع شهري، وإحصاء حسب الفئات \tabularnewline
Why (لماذا) & تدريب عملي + معرفة أين ذهبت الأموال \tabularnewline
How (كيف يُبنى) & Swift + SwiftUI + SwiftData + MVVM، محلي بالكامل دون خادم خلفي \tabularnewline
Constraint (القيود) & الإنجاز خلال أسبوع؛ الحد الأدنى للدعم iOS 17 (شرط SwiftData) \tabularnewline
\hline
\end{tabularx}
\end{center}

\begin{center}\small\color{gray}عناصر المتطلبات الخمسة لدفتر الحسابات\end{center}

ثم فككه إلى خطوات صغيرة «كل خطوة منها قابلة للاستخدام»: (1) بناء طبقة بيانات SwiftData (تستطيع تخزين عملية)؛ (2) صفحة القائمة (تستطيع عرض الكشف)؛ (3) نافذة الإدخال (تستطيع إضافة عملية)؛ (4) الحذف والتحرير؛ (5) بطاقة الملخص الشهري. وبعد كل خطوة يبقى التطبيق «شيئًا قابلًا للاستخدام».

\booksection{ثانيًا: طبقة البيانات — بيت لـ«الحساب»}

عرّف أولًا شكل العملية المالية الواحدة — فالحقول ينبغي أن تكفي دون طمع:

\begin{lstlisting}[language=Swift]
import SwiftData

@Model
class Record {
    var type: String      // "income" إيراد | "expense" مصروف
    var amount: Double    // المبلغ (يوان)
    var category: String  // التصنيف: طعام، مواصلات، راتب…
    var note: String      // ملاحظة (اختياري)
    var date: Date        // تاريخ التسجيل

    init(type: String, amount: Double, category: String, note: String = "", date: Date = Date()) {
        self.type = type
        self.amount = amount
        self.category = category
        self.note = note
        self.date = date
    }
}
\end{lstlisting}

\begin{tcbinfo}{نصيحة: لماذا نستخدم Double للمبلغ؟}
تستخدم تطبيقات المال الدقيقة النوع Decimal (عشري دقيق) لتجنب أخطاء الفاصلة العائمة (المشكلة القديمة 0.1+0.2$\neq$0.3). وفي مشروع تدريبي يكفي Double؛ وإذا جاء يوم تصنع فيه منتجًا ماليًا حقيقيًا فتذكر حينها التحول إلى Decimal أو «عدد صحيح Int بوحدة السنتات» — فخ حذرت منه أجيال من المبرمجين قبلك.

\end{tcbinfo}

\booksection{ثالثًا: MVVM — «فصل الواجهة عن المنطق»}

أن تجعل الواجهة تتعامل مع قاعدة البيانات مباشرة؟ يعمل، لكن كود الواجهة سيصير طويلًا فاحشًا. تجعل معمارية MVVM كلًا يؤدي دوره: الواجهة View ترسم فقط، ونموذج العرض ViewModel (مركز المنطق) يدير البيانات، والنموذج Model (SwiftData) يتولى التخزين.

العرض (واجهة SwiftUI)

ViewModel (مركز المنطق @Observable)

Model (طبقة بيانات SwiftData)

\begin{center}\small\color{gray}شكل 4: البنية ثلاثية الطبقات في MVVM — تتدفق البيانات في اتجاه واحد، ولا تجاوز أي طبقة حدودها\end{center}

\begin{lstlisting}[language=Swift]
import Observation

@Observable
class RecordViewModel {
    var records: [Record] = []
    private let context: ModelContext

    init(context: ModelContext) { self.context = context }

    func load() {
        // تصفية بـ #Predicate من SwiftData + ترتيب
        let descriptor = FetchDescriptor<Record>(
            sortBy: [SortDescriptor(\.date, order: .reverse)]
        )
        records = (try? context.fetch(descriptor)) ?? []
    }

    func add(type: String, amount: Double, category: String, note: String) {
        let r = Record(type: type, amount: amount, category: category, note: note)
        context.insert(r)
        try? context.save()
        load()                       // إعادة تحميل، وتتحدث الواجهة
    }

    func delete(_ record: Record) {
        context.delete(record)
        try? context.save()
        load()
    }
}
\end{lstlisting}

في الواجهة تكتب @State private var vm = RecordViewModel(...) ثم تعرض القائمة مباشرة من vm.records — فبمجرد تغير البيانات تنعكس الواجهة تلقائيًا.

\booksection{رابعًا: الواجهة — صفحة القائمة ونافذة الإدخال}

\begin{lstlisting}[language=Swift]
struct RecordListView: View {
    @State private var vm: RecordViewModel
    @State private var showAdd = false

    var body: some View {
        NavigationStack {
            List {
                Section("رصيد الشهر") {          // بطاقة الملخص
                    SummaryRow(records: vm.records)
                }
                Section("الحركات") {
                    ForEach(vm.records) { r in
                        RecordRow(record: r)
                    }
                    .onDelete { idx in
                        idx.forEach { vm.delete(vm.records[$0]) }
                    }
                }
            }
            .navigationTitle("دفتر مصروفاتي")
            .toolbar {
                Button { showAdd = true } label: {
                    Image(systemName: "plus")
                }
            }
            .sheet(isPresented: $showAdd) {
                AddRecordView(vm: vm)          // نافذة الإدخال
            }
        }
    }
}
\end{lstlisting}

نقاط نافذة الإدخال: لوحة مفاتيح حقل المبلغ من نوع .keyboardType(.decimalPad) (لوحة أرقام)، والدخل/المصروف عبر أداة مقسمة (Picker)، والفئة عبر قائمة اختيار، وعند الضغط على «حفظ» تستدعي vm.add(...) ثم dismiss() لإغلاق النافذة.

\begin{tcbwarn}{تنبيه: لا تكاسل في التحقق من المدخلات}
حقل المبلغ يجب التحقق منه حتمًا: اعترض القيم الفارغة والسالبة والمدخلات غير الصالحة مثل “abc” (بالحكم عبر Double(text) الذي يعيد nil عند الفشل). فمن أكثر أسباب انهيار التطبيقات شيوعًا «معاملة مدخلات المستخدم كطفل مطيع». وتذكر: كل ما يدخله المستخدم عاملُه كعدو يخضع للتفتيش.

\end{tcbwarn}

\booksection{خامسًا: التحرير والحذف والملخص}

\begin{itemize}
  \item الحذف: .onDelete في List للحذف بسحب لليسار + نافذة تأكيد (فالحذف الخاطئ أكثر ما يكره مستخدمو تطبيقات الحسابات)؛
  \item التحرير: النقر على صف يفتح نافذة تحرير معبأة مسبقًا بالقيم الأصلية (تصحيح حرف واحد أسهل من إعادة الإدخال)؛
  \item الملخص: حساب «دخل الشهر / مصروفه / الرصيد» عبر filter مع reduce وعرضها في بطاقة أعلى الصفحة:
\end{itemize}

\begin{lstlisting}[language=Swift]
func monthSummary(_ records: [Record]) -> (income: Double, expense: Double) {
    let cal = Calendar.current
    let monthRecords = records.filter {
        cal.isDate($0.date, equalTo: Date(), toGranularity: .month)
    }
    let income = monthRecords
        .filter { $0.type == "income" }
        .reduce(0) { $0 + $1.amount }
    let expense = monthRecords
        .filter { $0.type == "expense" }
        .reduce(0) { $0 + $1.amount }
    return (income, expense)
}
\end{lstlisting}

\booksection{سادسًا: القبول — فحص بثلاث طبقات وتحديان متقدمان}

بعد إنهاء المشروع اقبله على ثلاث طبقات وفق «الوظائف الحدود المعمارية»: هل يسجل وهل يحذف وهل تبقى البيانات بعد إعادة تشغيل التطبيق (الوظائف)؛ هل ينهار عند إدخال 0 أو رقم سالب في المبلغ (الحدود)؛ هل اختلطت الواجهة بالبيانات وهل الحقول مكتوبة بشكل ثابت داخل الكود (المعمارية).

بعد نجاح النسخة الأساسية في القبول، أضف لنفسك «طبقين متقدمين»:

\begin{enumerate}
  \item إحصاء الفئات: أضف علامة تبويب «إحصاءات»، وارسم مصروفات الفئات دائريًا أو أعمدة باستخدام Swift Charts — مكتبة الرسوم الرسمية من Apple، بضعة أسطر كود تنتج الرسم؛
  \item مزامنة iCloud: تفعيل مزامنة iCloud في SwiftData لا يتطلب سوى ضبط بضع إعدادات — فيتزامن الدفتر تلقائيًا بين iPhone وiPad، ويرتقي المشروع فورًا إلى مستوى «راقٍ».
\end{enumerate}

\begin{tcbpractice}{هيا نجرّب: أكمله خلال هذا الأسبوع}
انطلق بالترتيب «طبقة البيانات ViewModel صفحة القائمة نافذة الإدخال بطاقة الملخص»، خطوة كل يوم. وإذا علقت فاطرح رسالة الخطأ كما هي على محرك البحث أو على الذكاء الاصطناعي (أتذكر الجزء السابع «البرمجة العملية بالذكاء الاصطناعي»؟ عند وصف المشكلة وضح «السلوك المتوقع، والسلوك الفعلي، وخطوات إعادة الإنتاج»). وبعد الانتهاء سجل به حسابات أسبوع كامل — فاستخدام تطبيق صنعته بنفسك يولد شعور الإنجاز الذي يدفعك إلى مواصلة الطريق.

\end{tcbpractice}

\begin{tcbsummary}{خلاصة الفصل}
\begin{itemize}
  \item عناصر المتطلبات الخمسة تفكيكها إلى خطوات صغيرة «كل خطوة قابلة للاستخدام» ثم الشروع في العمل
  \item طبقة البيانات: نموذج @Model + حاوية modelContainer + استعلام FetchDescriptor
  \item توزيع أدوار MVVM: الواجهة View ترسم فقط، وViewModel المبني على @Observable يدير المنطق، وSwiftData يتولى التخزين
  \item الواجهة: List مع Section مع toolbar ونافذة الإدخال sheet
  \item قبول بثلاث طبقات (الوظائف / الحدود / المعمارية) + تحديان متقدمان (Swift Charts ومزامنة iCloud)
\end{itemize}

\end{tcbsummary}
